% Figure from MATH 493 notes, section 52. Compile with the notes preamble.
\begin{tikzpicture}[font=\small]
  \node[draw=figR, very thick, rounded corners=3pt, fill=figR!8, inner sep=4pt] (P) at (0,0) {$P = \langle (1\,2\,3) \rangle$};
  \node[figR, font=\scriptsize, align=center] at (0,-0.75) {fixed by $P$:\\ orbit of size $1$};
  \node[draw=figB, rounded corners=3pt, fill=figB!6, inner sep=3pt] (A) at (4.2,1.1) {$\langle (1\,2\,4) \rangle$};
  \node[draw=figB, rounded corners=3pt, fill=figB!6, inner sep=3pt] (B) at (6.2,-0.9) {$\langle (2\,3\,4) \rangle$};
  \node[draw=figB, rounded corners=3pt, fill=figB!6, inner sep=3pt] (C) at (2.6,-0.9) {$\langle (1\,3\,4) \rangle$};
  \draw[->, figB, thick, bend left=25] (A) to node[right, font=\scriptsize, text=black] {$c_{(1\,2\,3)}$} (B);
  \draw[->, figB, thick, bend left=25] (B) to (C);
  \draw[->, figB, thick, bend left=25] (C) to (A);
  \node[figB, font=\scriptsize] at (4.4,-1.75) {one orbit of size $3$};
  \draw[black!30, rounded corners=8pt] (-1.6,-1.3) rectangle (7.6,1.75);
  \node[black!60, font=\scriptsize, anchor=north west] at (-1.55,1.7) {$X$ = Sylow $3$-subgroups of $S_4$};
  \node[font=\scriptsize] at (3,-2.45) {$|X| = 1 + 3 = 4 \equiv 1 \pmod 3$};
\end{tikzpicture}
